\hypertarget{conclusion}{%
\chapter{Conclusion}\label{conclusion}}

This dissertation proposed and evaluated IDRISK2, an end-to-end pipeline for the automated assignment of FoodEx2 classification codes to food images, combining optical character recognition, vision-language reasoning, and a retrieval-augmented generation layer grounded in the official EFSA FoodEx2 taxonomy. This concluding chapter summarises the research carried out, the contributions it makes to the field, and the closing reflections that arise from the evaluation.

\hypertarget{summary-of-research}{%
\section{Summary of Research}\label{summary-of-research}}

The research was motivated by the persistent difficulty of mapping food product information to FoodEx2 codes in a regulatory inspection context. Manual classification is time-consuming, error-prone, and creates a bottleneck for downstream risk assessment by EFSA. Recent advances in vision-language and retrieval-augmented architectures opened the possibility of automating this mapping, but no existing system had been demonstrated end-to-end from real-world product imagery to compliant FoodEx2 outputs within a closed deployment environment aligned with GDPR and EU AI Act requirements. The central research question formulated in Section 1.3 asked how the integration of computer vision, OCR, and retrieval-augmented language models could improve the efficiency, accuracy, and security of FoodEx2 coding while remaining compliant with the cybersecurity requirements of ASAE, SIT, and FVST.

To address that question, Chapter 4 specified a modular pipeline comprising five processing layers: image preprocessing and OCR (Sections 4.2 and 4.3), vision-language interpretation (Section 4.4), Advanced RAG retrieval over a knowledge base built from the EFSA technical PDFs and the Appendix B taxonomy XLSX (Sections 4.5.1 and 4.5.2), structured FoodEx2 generation with audit logging (Section 4.5.3), and a closed security and compliance layer enforcing role-based access control, data minimisation, and tamper-evident audit chaining (Section 4.6). The implementation was evaluated empirically on 34 photographs of culturally specific dishes and packaged products, against a manually annotated FoodEx2 ground truth, and the results were reported in Chapter 5.

The evaluation established four substantive findings. First, the Vision-Language Model produces preliminary FoodEx2 codes that are structurally well-formed but semantically fabricated in 100\% of evaluations, with hallucinated identifiers presented at high confidence. Second, the RAG retrieval and reranking layers correct these hallucinations and contribute the entirety of the system's useful classification accuracy, achieving 29.4\% exact-match and 58.8\% family-level accuracy. Third, the system's confidence signal does not track actual correctness, with a false high-confidence rate of 71.9\% and a silent error rate of 95.8\% that materially weakens the human-oversight mechanism. Fourth, the closed deployment architecture, the data-minimisation pipeline, and the cryptographically chained audit log satisfy the explainability and traceability requirements of GDPR Article 22 and the documentation obligations of the EU AI Act for high-risk AI systems, even where the system's accuracy and calibration fall short of fully autonomous deployment.

\hypertarget{contributions-to-the-field}{%
\section{Contributions to the Field}\label{contributions-to-the-field}}

Relative to the four gaps identified in Section 3.5, this work makes the following contributions.

First, this dissertation presents the first end-to-end pipeline reported in the literature that ingests raw food product imagery, as captured under real-world inspection conditions, and produces structured FoodEx2 classifications grounded in the official EFSA taxonomy. The most advanced prior system, FEAST \parencite{lorenzo_molfetta_47d7c4cd}, operates on pre-structured textual food descriptions and does not address the upstream visual-interpretation problem; the IDRISK2 system covers the entire stack from photograph to compliant code assignment.

Second, this dissertation contributes the first systematic empirical evaluation of Vision-Language Model behaviour on FoodEx2 classification, spanning the proprietary-versus-open-source dimension and covering both the base term and the facet dimensions of the compound classification. The finding that GPT-4o, the strongest available proprietary VLM at the time of writing, achieves zero exact-match accuracy and emits hallucinated taxonomy codes in 100\% of cases (94\% misattributed real codes plus 6\% non-existent codes), and that this zero-accuracy result reproduces under Claude Sonnet 4.6 and under the LLaVA-NeXT-Mistral-7B open-source baseline, elevates the hallucination phenomenon from a single-model observation to a class-level property of contemporary Vision-Language Models on structured-taxonomy classification tasks. The complementary finding that the RAG layer compensates for this VLM failure to recover 29.4\% exact-match accuracy under the proprietary VLMs but only 2.9\% under LLaVA-NeXT-7B delimits the empirical boundary at which proprietary-to-open-source VLM substitutability breaks down for this task, at approximately one order of magnitude of system accuracy. The facet-level evaluation reported in Section 5.7.4 extends this characterisation to the descriptive dimension of the FoodEx2 code, documenting that the same proprietary-versus-open-source ordering reproduces (per-image facet hit rate of 38.2\% for GPT-4o, 35.3\% for Claude, 8.8\% for LLaVA-NeXT-7B) and identifying two architectural defects of independent interest, the systematic non-emission of the packaging-format and packaging-material facets (0.0\% coverage across all three VLMs even on relevant packaged-drink inputs), and the high coverage but low accuracy of the ingredient facet (77--94\% emission rate at 12--20\% accuracy when emitted), which characterises a schema-validity-without-semantic-correctness regime extending from the base term to the facet dimension of the classification. This characterisation is novel, methodologically replicable, and directly motivates the use of retrieval-augmented architectures, evaluated against the strongest VLM available, for any structured-taxonomy classification task in this domain.

Third, this dissertation operationalises the regulatory compliance dimensions that the literature on automated food classification has largely overlooked. The data-minimisation record schema, the cryptographically chained audit log, and the role-based access control framework collectively demonstrate that the security, logging, and access-control infrastructure required to support FoodEx2 classification in a regulatory context can be implemented as a production-grade architectural layer, addressing the foundational data-residency, audit-traceability, and access-control obligations of GDPR Article 22 and aligning with the human-oversight, transparency, and documentation principles of EU AI Act Articles 14 and 13. The qualifier production-grade applies to this compliance and security infrastructure rather than to the classification engine itself, whose accuracy and calibration (Sections 5.6 and 6.4) place it within the decision-support rather than the autonomous-classifier regime. The contribution should therefore be read as a foundational step toward full high-risk AI compliance, the completion of which would require an extensive external conformity assessment, alignment with the relevant ISO standards (notably ISO/IEC 42001), and a quality management system that extends substantially beyond the architectural components implemented in the present prototype.

Fourth, this dissertation specifies and implements a closed, self-hosted deployment architecture that operates entirely on infrastructure controlled by the deploying authority, without reliance on public cloud services or external unverified APIs for the storage, indexing, or processing of product data. The closed posture, Chroma vector store on local disk, taxonomy XLSX and EFSA PDFs ingested at deployment time, VLM and LLM API calls restricted to the model providers' inference endpoints with no persistence of product data outside the deploying authority, addresses the data-governance constraints that disqualify cloud-dependent alternatives from official regulatory deployment in ASAE, SIT, and FVST contexts.

In addition to the four primary contributions aligned with the gaps of Section 3.5, the dissertation surfaces four findings of independent methodological value. The first is a heterogeneous-chunking knowledge-base construction methodology that combines semantically chunked regulatory documents with atomic taxonomy-term nodes indexed at term granularity; the empirical evaluation demonstrates that this design choice is essential to retrieval effectiveness, with every reranked candidate being a taxonomy term and no documentation chunk ever surfacing above the term nodes. The second is the documentation of the label-leak architectural defect as a quantitatively measurable failure mode in retrieval-augmented structured-output systems. The third is the empirical characterisation of the cultural-coverage limitations of the FoodEx2 taxonomy with respect to Lusophone cuisine. The fourth is the demonstration that a confidence-based human-review trigger drawn solely from the VLM's self-reported confidence is insufficient to satisfy the human-oversight obligations of EU AI Act Article 14 in practice.

\hypertarget{final-remarks}{%
\section{Final Remarks}\label{final-remarks}}

The IDRISK2 system, as evaluated in this dissertation, is not ready for autonomous deployment in a regulatory inspection workflow. An exact-match accuracy of 29.4\% and a silent error rate of 95.8\% place the system below the operational threshold at which inspectors could rely on its classifications without review. However, the same evaluation establishes that the system's family-level accuracy of 58.8\% combined with its mandatory audit trail produces real operational value when the system is deployed as a decision-support tool with human review on every output: it reduces the inspector's effective search space from 31,303 taxonomy terms to a small reranked candidate set, it produces a complete provenance record for each candidate, and it does so within a closed, GDPR-compliant infrastructure at a median latency of 22 seconds per image.

The dissertation's contribution is therefore best read not as a claim that automated FoodEx2 classification is solved, but as a claim that the architecture required to make automated FoodEx2 classification operationally viable can now be specified with empirical precision. The retrieval-augmented design is necessary; the closed deployment is feasible; the regulatory compliance is achievable; the accuracy and calibration remain bound to the limitations of the available vision-language models and the cultural coverage of the underlying taxonomy. Each of these residual limitations admits to a concrete future-work programme, which Chapter 8 develops in turn.

Finally, the dissertation closes by returning to the broader motivation articulated in Section 1.2: the goal of reducing the inspector's manual workload, of improving the harmonisation of food classification data collected by the authorities that report to EFSA, and of advancing the application of multimodal AI to regulatory food safety. The IDRISK2 system contributes a working architectural proposal for that goal, an honest empirical accounting of its present limitations, and a documented programme of future development that would close the gap between its current decision-support capability and the autonomous classifier that the regulatory operational context will eventually require.
